\section{Tipo de Pentest}

\paragraph{}
Considerando as informações obtidas durante a fase de reconhecimento e os
resultados da varredura de vulnerabilidades realizada com o OWASP ZAP, o tipo
de teste de penetração mais adequado para o Laboratório de Segurança é o
\textit{Gray Box Pentest} com foco em aplicações web.

\paragraph{}
Essa abordagem é adequada porque o analista possui conhecimento parcial do
ambiente antes da execução dos testes, incluindo informações sobre a aplicação
alvo, URLs acessíveis, parâmetros utilizados, funcionalidades disponíveis e
vulnerabilidades previamente identificadas. Ao mesmo tempo, ainda existe a
necessidade de realizar atividades de enumeração, validação e exploração
controlada para confirmar os riscos encontrados.

\paragraph{}
Os resultados da análise inicial identificaram vulnerabilidades críticas e de
alta prioridade, como \textit{SQL Injection}, \textit{Cross Site Scripting
(Reflected)}, \textit{Path Traversal} e \textit{External Redirect}, indicando
que a principal superfície de ataque encontra-se na aplicação web
\texttt{http://testasp.vulnweb.com}. Dessa forma, o escopo do Pentest deve
priorizar testes de segurança em aplicações web (\textit{Web Application
Penetration Testing}), com foco na validação da exploração dessas falhas e na
avaliação de seus impactos sobre a confidencialidade, integridade e
disponibilidade das informações.

\paragraph{}
Não foram identificados elementos suficientes que justifiquem a realização de
testes específicos de segurança móvel ou avaliações aprofundadas de
infraestrutura de rede. Portanto, os esforços serão concentrados na análise da
camada de aplicação, onde foram observados os principais riscos e
vulnerabilidades do ambiente avaliado.

\paragraph{}
A adoção de um Pentest do tipo \textit{Gray Box} permite otimizar os recursos e
o tempo disponíveis, direcionando as atividades para os componentes mais
críticos do ambiente e aumentando a eficiência na identificação e validação das
vulnerabilidades previamente detectadas.
